\appendix

\section{Duality and partial trace}
\label{sec:duality}

In this appendix we show that given a kernel $\Phi\colon X\ra Y$ and
its dual kernel \(\Phi^\chk\colon Y \ra X\) there are canonical
2-morphisms
\[ \SerK_X^{-1} \ra \Phi^\chk \circ \Phi\quad\text{and}\quad \Phi
\circ \SerK_X \circ \Phi^\chk \ra \Id_Y \]
giving rise to a variety of natural adjunctions satisfying
a number of compatibility relations.

The notion of duality in $\gVar$ is seen to be a middle-ground between
the operations $\tau_L$ and $\tau_R$: it is an involution, unlike
$\tau_L$ and $\tau_R$, but it does not respect composition, which
$\tau_L$ and $\tau_R$ do.

 \subsection{Polite duality.}
Recall from Section~\ref{subsec:serXY} that for every space $X$ there
is the Serre kernel $\SerK_X\colon X \ra X$ such that for spaces $X$
and $Y$ the functor
\[ \SerK_Y \circ {-} \circ \SerK_X\colon \bHom_\gVar(X, Y) \ra
\bHom_\gVar(X, Y) \]
is a Serre functor for the category $\bHom_\gVar(X,Y)$.

\begin{Definition}
\label{def:dualpair}
If $\Phi\colon X \ra Y$ and $ \pd{\Phi}\colon Y \ra X$ are kernels then
a \emph{polite duality} between them, denoted \( \Phi \dual \pd{\Phi}
\), consists of adjunctions as follows (numbered as shown)
\begin{gather*}
  \pd{\Phi} \circ \SerK_Y \adjoint_1 \Phi \adjoint_2 \SerK_X \circ
  \pd{\Phi},
  \\
  \Phi \circ \SerK_X \adjoint_3 \pd{\Phi} \adjoint_4 \SerK_Y \circ \Phi,
\end{gather*}
such that the following compatibility relations hold.
\begin{quote}
\begin{itemize}
\item[(1+2).] For kernels $\Theta\colon Z\ra X$ and $\Psi\colon Z\ra
  Y$, the diagram of isomorphisms below commutes:
  \[
  \begin{diagram}[height=2em,width=2em,labelstyle=\scriptstyle]
    \Hom(\Psi, \Phi \circ\Theta) & \rTeXto^1 &
    \Hom(\pd{\Phi} \circ\SerK_Y\circ \Psi, \Theta) \\
    \dTeXto^{\text{Serre}} & &
    \dTeXto^{\text{Serre}} \\
    \Hom(\Phi\circ\Theta, \SerK_Y\circ \Psi \circ \SerK_Z)^\chk &
    \rTeXto^2 & \Hom(\Theta, \SerK_X\circ \pd{\Phi}\circ \SerK_Y\circ
    \Psi\circ \SerK_Z)^\chk.
  \end{diagram}
  \]
\item[(2+3).] The composite map
  \[ \Hom(\Theta, \pd{\Phi}\circ\Psi) \xrightarrow[3]{\sim}
  \Hom(\Phi\circ\SerK_X\circ\Theta, \Psi) \xrightarrow[2]{\sim}
  \Hom(\SerK_X\circ\Theta, \SerK_X\circ \pd{\Phi}\circ \Psi) \]
  is the one induced by composition with $\SerK_X$.
\item[(3+4).] Same as (1+2), but for adjunctions 3 and 4.
\item[(1+4).] Same as (2+3), but for adjunctions 1 and 4.
\end{itemize}
\end{quote}
\end{Definition}


It is useful to think of this definition in terms of string diagrams.
For each of the four adjunctions we get a unit and a counit.  Denoting
$\Phi$ by an upward oriented line and $\pd{\Phi}$ by a downward oriented
line, we can draw the units and counits as follows.
\begin{gather*}
\eta_1\colon\pic{eta1}\quad
\eta_2\colon\pic{eta2}\quad
\eta_3\colon\pic{eta3}\quad
\eta_4\colon\pic{eta4}\\
\epsilon_1\colon\pic{epsilon1}\quad
\epsilon_2\colon\pic{epsilon2}\quad
\epsilon_3\colon\pic{epsilon3}\quad
\epsilon_4\colon\pic{epsilon4}
\end{gather*}

Relation (2+3) can be represented graphically in the following way:
for any $\alpha\in\Hom(\Theta,\pd{\Phi}\circ\Psi)$ we have
\[\pic{R2and3a}=\pic{R2and3b}.\]
From this we can deduce the equality
\[\pic{Y23a}=\pic{Y23b},\]
so we define the following diagram to be this common morphism:
\[\gamma_\Phi:\qquad\pic{Y23c}.\]
Similarly we can deduce that the counits $\epsilon_2$ and $\epsilon_3$
are equal, and we define the following diagram to be the common
morphism:
\[\rote_\Phi:\qquad\pic{M23}.\]
Likewise, using the relation (1+4), from the units $\eta_1$ and
$\eta_4$, and the counits $\epsilon_1$ and $\epsilon_4$ we obtain the
common morphisms denoted as follows.
\[\gamma_{\Phi^\chk}:\quad\pic{Y14}\qquad \rote_{\Phi^\chk}:\quad \pic{M14}\]
So the four units and four counits are obtained from these two
M-shaped and two Y-shaped morphisms.

Relations (1+2) and (3+4) in the definition of polite duality are
essentially equivalent to the invariance of the Serre trace under a
partial trace in the following sense.  Given a polite duality
$\Phi\dual\pd{\Phi}$ and a morphism $\alpha\in
\Hom(\Phi\circ\Theta,\SerK\circ \Phi\circ \Psi)$ we can define the
left partial trace in $\Hom(\Theta,\SerK\circ \Psi)$ with respect to
$\pd{\Phi}$ as drawn below.
\[\pic{PartialTraceLeft}\]
Similarly we can define a right partial trace when the $\Phi$ is on
the right rather than on the left.  So the following diagram is the
right partial trace with respect to $\Phi$ and $\pd{\Phi}$ of a morphism
$\alpha'\in \Hom(\Theta'\circ\Phi,\Psi'\circ \Phi\circ\SerK)$.
\[\pic{PartialTraceRight}\]

\begin{Theorem}
  If $\Phi$ and $\pd{\Phi}$ form a politely dual pair of kernels, then
  the Serre trace is invariant under partial trace with respect to
  $\Phi$:
  \[\Tr\left(\pic{ApPT1}\right)=\Tr\left(\pic{ApPT2}\right);
  \quad \Tr\left(\pic{ApPT1b}\right)=\Tr\left(\pic{ApPT2b}\right).\]
\end{Theorem}

\begin{Proof}
  We will just prove the left partial trace case.  Relation (1+2) says
  that for $\beta\in \Hom (\Psi,\Phi\circ\Theta)$ and $\gamma\in
  (\Theta,\SerK\circ\pd{\Phi}\circ\SerK\circ\Psi\circ\SerK)$ we have the
  equality
  \[\Tr\left(\pic{ApPT3}\right)=\Tr\left(\pic{ApPT4}\right).\]
  Applying this with
  \[\beta:=\pic{ApPT6}\quad\text{and}\quad
  \gamma:=\pic{ApPT5}\] we obtain the equality
  \[\Tr\left(\pic{ApPT7}\right)=\Tr\left(\pic{ApPT8}\right).\]
  The result follows from the commutativity property of the Serre
  trace, namely that $\Tr(\beta\circv \alpha)=\Tr(\SerF(\alpha)\circv
  \beta)$.
\end{Proof}

\begin{Proposition}
\label{prop:compduals}
  Polite dualities $\Phi \dual \pd{\Phi}$ and $\Psi \dual \pd{\Psi}$ for
  two kernels $\Phi\colon X\ra Y$ and $\Psi\colon Y\ra Z$ canonically
  induce a polite duality
  \[ \Psi \circ \Phi \dual \pd{\Phi} \circ \SerK_Y \circ \pd{\Psi}. \]
\end{Proposition}

\begin{Proof}
  The required adjunctions are constructed in the obvious fashion,
  using the fact that the composition, in reverse order, of adjoints
  (left or right) of composable functors is naturally an adjoint (in
  the same direction) of the composition of the functors.  The
  compatibilities required by the polite duality follow from straight
  forward checks.
\end{Proof}

Adjunctions are intrinsic parts of a polite duality.  Conversely, the
following proposition shows that any adjunction induces a polite
duality in a natural way.

\begin{Proposition}
\label{prop:adjgivesdual}
Given an adjunction $\Psi \adjoint \Phi$, define $\pd{\Phi}:= \Psi \circ
\SerK_Y^{-1}$.  Then there exists a polite duality $\Phi \dual
\pd{\Phi}$ where the adjunction $\adjoint_1$ is the given one
$\Psi\adjoint \Phi$.
\end{Proposition}

\begin{Proof}
  As $\SerK_X\circ \pd{\Phi}= \SerK_X\circ\Psi\circ\SerK^{-1}_Y$, compatibility relation (1+2) yields an adjunction
  \[ \Phi \adjoint_2 \SerK_X
  \circ \pd{\Phi}. \] Similarly, relations (1+4) and (2+3) force
  adjunctions
  \[ \Phi \circ \SerK_X \adjoint_3 \pd{\Phi} \adjoint_4 \SerK_Y \circ
  \Phi, \]
  and it is easy to see that relation (3+4) is then automatically
  satisfied.
\end{Proof}

\begin{Proposition}
\label{prop:existdual}
Let $\Phi\in\D(X\times Y)$ be a kernel, and let
$\Phi^\chk\in\D(Y\times X)$ be the dual $\mathbf{R}\sHom_{X\times
  Y}(\Phi, \cO_{X\times Y})$, regarded as a kernel from $Y$ to $X$.
Then there exists a canonical polite duality \(\Phi \dual \Phi^\chk\).
\end{Proposition}

\begin{Proof}
For kernels $\Theta \in \D(Z\times X)$ and $\Psi\in \D(Z\times Y)$
consider the sequence of isomorphisms
\begin{align*}
\Hom_{\D(Z\times Y)} & (\Psi, \Phi \circ \Theta) \\*
  & \iso \Hom_{\D(Z\times Y)}(\Psi, \pi_{ZY, *}
      (\pi_{ZX}^* \Theta \otimes \pi_{XY}^* \Phi)) \\
  & \iso \Hom_{\D(Z\times X)}(\pi_{ZX,!}
    (\pi_{ZY}^* \Psi \otimes \pi_{YX}^* \Phi^\chk), \Theta) \\
  & = \Hom_{\D(Z\times X)}(\pi_{ZX, *}
      (\pi_{ZY}^*\Psi \otimes \pi_{YX}^*
        (\Phi^\chk \otimes \omega_Y[\dim Y])), \Theta) \\
  & = \Hom_{\D(Z\times Y)}((\Phi^\chk\circ \SerK_Y)\circ\Psi, \Theta).
\end{align*}
Taking $Z = Y$, $\Theta = \Phi^\chk \circ \SerK_Y$ and $\Psi=\Id_Y$
yields a morphism {\em of kernels}
\[ \Id_Y \ra \Phi \circ \Phi^\chk\circ \SerK_Y; \]
in a similar fashion we obtain a morphism
\[ \Phi^\chk \circ \SerK_Y \circ \Phi \ra \Id_X. \]
These two morphisms satisfy the identities needed to make $\Phi^\chk
\circ \SerK_Y$ the left adjoint of $\Phi$.
Proposition~\ref{prop:adjgivesdual} gives the result.
\end{Proof}

 \subsection{Reflexively polite kernels.}
We still need to address one more compatibility between the dualities
constructed above.  Given a kernel $\Phi$, the previous proposition
yields a polite duality
\[ \Phi \dual \Phi^\chk. \]
Given any polite duality $\Phi\dual \pd{\Phi}$, we get, symmetrically, a polite
duality $\pd{\Phi}
\dual \Phi$ by switching the adjunctions $1$ and  $3$ and the adjunctions
$2$ and $4$.  Thus there is a natural polite duality
\[ \Phi^\chk \dual \Phi. \]
On the other hand, applying Proposition~\ref{prop:existdual} to the
kernel $\Phi^\chk$ we get a polite duality
\( \Phi^\chk \dual \Phi^{\chk\chk}\)
and then, using the canonical identification $\Phi^{\chk\chk} \iso
\Phi$, we get {\em another} polite duality
\[ \Phi^\chk \dual \Phi.\]
The fundamental question is whether these two polite dualities are the same.

\begin{Definition}
  Let $\Phi\colon X\ra Y$ be a kernel.  We shall say that $\Phi$ is
  \emph{reflexively polite} if the two dualities above are equal.
\end{Definition}

\noindent Immediately we get the following result.

\begin{Proposition}
  If a kernel $\Phi$ is reflexively polite then so is $\Phi^\chk$.
\end{Proposition}

\begin{Proposition}
\label{prop:politecompduals}
If $\Phi\colon X \ra Y$ and $\Psi\colon Y \ra Z$ are reflexively
polite kernels then so is $\Psi \circ \Phi$.
\end{Proposition}

\begin{Proof}
  This follows at once from the fact that the adjunctions defined in
  Proposition~\ref{prop:compduals} are canonical.
\end{Proof}

\begin{Proposition}
\label{prop:dualbox}
Suppose that $\Phi\colon X\ra Y$ and $\Psi\colon X'\ra Y'$ are
reflexively polite kernels then so is $\Phi\boxtimes \Psi\colon X
\times X' \ra Y\times Y'$, the kernel defined by
\[ \Phi\boxtimes \Psi := \pi_{XY}^* \Phi \otimes \pi_{X'Y'}^*\Psi. \]
\end{Proposition}

\begin{Proof}
  This is obvious once one realises that all operations decompose with
  respect to the box product operation, and the canonical bundle of a
  product of spaces is the box product of the canonical bundles of the
  factors.
\end{Proof}

\begin{Proposition}
\label{prop:dualker}
Let $\Delta\colon X \ra X\times X$ denote the diagonal map, let $\cE$
be an object in $\D(X)$, then the kernel $\Delta_*\cE\colon X \ra X$
is reflexively polite.
\end{Proposition}

\begin{Proof}
  The dual $(\Delta_*\cE)^\chk$ of $\Delta_*\cE$ is naturally
  identified with 
  \[ \Delta_* (\cE^\chk \otimes \omega_X^{-1}[-\dim X]). \] The result
  now follows from Proposition~\ref{prop:deltalstar} and the fact that
  the functor given by tensoring with an object of the derived
  category is left and right adjoint to the one obtained by tensoring
  with the dual object.
\end{Proof}

\begin{Proposition}
\label{prop:dualcox}
Let $X$ be a space, then the kernel $\cO_X$ considered as a kernel
$X\ra \pt$ is reflexively polite, and so also is $\cO_X^\chk\colon
\pt\to X$.
\end{Proposition}

\begin{Proof}
  First, observe that checking reflexive duality amounts to checking
  the equality of two quadruples of adjunctions.  In each quadruple,
  any one of the adjunctions determines the remaining three, and thus
  in order to check the equality it suffices to check that one of the
  compatibilities (1+2), (1+3), {etc}.\ holds with one adjunction
  chosen from one duality, and the other one from the second duality.

  Furthermore, adjunctions are completely determined by their
  respective units and counits, and these are determined by considered
  functors {\em between derived categories} of the form
  \[ \Phi \circ {-} \colon \D(Z\times X) =:\bHom_\gVar(Z, X) \ra
  \bHom_\gVar(Z, \pt) := \D(Z), \]
  and analogues for $\Phi^\chk \circ {-}$, for various choices of the
  space $Z$.  Thus, if we argue that the desired equality holds for
  the adjunctions between these induced functors (for arbitrary choice
  of $Z$), we will have argued that $\Phi$ is reflexively polite.

  For ease of notation, write $\Phi$ for the kernel $\cO_X\colon X\to \pt$.  For a given space $Z$, let $\pi_Z\colon X\times Z \ra Z$ denote the
  projection.  The functors $\Phi \circ {-}$ and $\Phi^\chk \circ
  {-}$ are then naturally identified with $\pi_{Z,*}$ and $\pi_Z^*$,
  respectively.  Adjunctions $1$ and $2$ from the polite duality
  $\Phi\dual \Phi^\chk$ correspond to the classical adjunctions
  \begin{align*}
    \pi_Z^* & \adjoint \pi_{Z,*} \adjoint \pi_Z^!, \intertext{while
      the same ones from the polite duality $\Phi^\chk\dual
      \Phi^{\chk\chk}$ correspond to} \pi_{Z,!} & \adjoint \pi_Z^*
    \adjoint \pi_{Z,*}.
  \end{align*}
  Indeed, the standard way (see~\cite[Theorem 4.6]{CalSkim}) to define
  the adjunctions $\pi_{Z,*} \adjoint \pi_Z^!$ and $\pi_{Z,!}\adjoint
  \pi_Z^*$ is to require them to satisfy the analogue of condition
  (1+2) from the definition of a polite duality.

  The condition of polite duality can now can be stated as the
  statement that the composite isomorphism
  \begin{align*}
    \Hom_{X\times Z}(\Theta, & \pi_Z^* \Psi) \iso \Hom_Z(\pi_{Z,!}
    \Theta, \Psi) =
    \Hom_Z\left(\pi_{Z,*}\SerF(\Theta), \Psi\right) \\
    & \iso \Hom_{X\times Z}\left(\SerF(\Theta), \pi_Z^! \Psi\right)
    = \Hom_{X\times Z}\left(\SerF(\Theta), \SerF (\pi_Z^*
      \Psi)\right)
  \end{align*}
  is the one induced by the functor $\SerF({-}) = {-} \otimes
  \pi_X^* \omega_X[\dim X]$, where $\pi_X$ denotes the projection from
  $X\times Z$ to $X$.  This fact corresponds to the fact that the
  diagram marked (!) below commutes {\small
    \[
    \begin{diagram}[height=2em,width=2em,labelstyle=\scriptstyle]
      \Hom(\Theta, \pi^* \Psi) & \rLine^\sim & \Hom(\SerF ^{-1} \pi_*
      \SerF \Theta, \Psi) &
      \rLine^\sim & \Hom(\pi_* \SerF \Theta, \SerF \Psi) & \rLine^\sim & \Hom(\SerF \Theta, \SerF \pi_* \Psi) \\
      \dLine^\sim & \mbox{Serre} & \dLine^\sim & \mbox{(!)} &
      \dLine^\sim &
      \mbox{Serre} & \dLine^\sim \\
      \Hom(\pi^*\Psi, \SerF \Theta)^\chk & \rLine^\sim & \Hom(\Psi,
      \pi_*\SerF \Theta)^\chk & \rEq & \Hom(\Psi, \pi_*\SerF
      \Theta)^\chk & \rLine^\sim & \Hom(\pi^*\Psi, \Theta)^\chk.
    \end{diagram}
    \]}

\end{Proof}

\begin{Proposition}
\label{prop:pulldualpush}
Let $\pi_X\colon X \times Y \ra X$ be the projection, and (abusing
notation) denote by $\pi_{X,*}\colon X\times Y \ra X$ and
$\pi_X^*\colon X\ra X\times Y$ the kernels represented in $\D(X\times
X \times Y)$ by the structure sheaf of the graph of $\pi_X$.  Then
$\pi_{X,*}$ and $\pi_X^*$ are reflexively polite.
\end{Proposition}

\begin{Proof}
  Both kernels are of the form $\cO_{\Delta_X} \boxtimes \cO_Y$, and
  the result then follows from
  Propositions~\ref{prop:dualcox},~\ref{prop:dualker},
  and~\ref{prop:dualbox}.
\end{Proof}

\begin{Theorem}
  Every kernel is reflexively polite.
\end{Theorem}

\begin{Proof}
  With notation as in Proposition~\ref{prop:pulldualpush}, any kernel
  $\Phi\colon X\to Y$ decomposes as
\[ \Phi = \pi_{Y,*} \circ ({-} \otimes \Phi) \circ \pi_X^*. \]
Since the claim of the theorem holds for each individual kernel in the
decomposition (Propositions~\ref{prop:pulldualpush}
and~\ref{prop:dualker}), the result follows from
Proposition~\ref{prop:compduals}.
\end{Proof} 